\documentclass[11pt]{article}
\usepackage[a4paper, margin=2.6cm]{geometry}
\usepackage{amsmath, amssymb, bm}
\usepackage{microtype}
\usepackage{xcolor}
\usepackage[colorlinks=true, linkcolor=blue!50!black]{hyperref}

% -- notation: bold upright for matrices, bold italic for vectors -------------
\newcommand{\mA}{\bm{A}}
\newcommand{\mAo}{\bm{A}_0}
\newcommand{\mX}{\bm{X}}
\newcommand{\mS}{\bm{S}}
\newcommand{\mG}{\bm{G}}
\newcommand{\mI}{\bm{I}}
\newcommand{\mP}{\bm{P}}
\newcommand{\mPp}{\bm{P}_{\!\perp}}
\newcommand{\mZ}{\bm{Z}}
\newcommand{\mU}{\bm{U}}
\newcommand{\mV}{\bm{V}}
\newcommand{\mW}{\bm{W}}
\newcommand{\mY}{\bm{Y}}
\newcommand{\mSig}{\bm{\Sigma}}
\newcommand{\mM}{\bm{M}}
\newcommand{\mT}{\bm{T}}
\newcommand{\mLam}{\bm{\Lambda}}
\newcommand{\mDel}{\bm{\Delta}}
\newcommand{\pool}{\bm{\Pi}}
\newcommand{\liftop}{\bm{E}}
\newcommand{\vx}{\bm{x}}
\newcommand{\vv}{\bm{v}}
\newcommand{\vc}{\bm{c}}
\newcommand{\ve}{\bm{e}}
\newcommand{\R}{\mathbb{R}}
\newcommand{\norm}[1]{\lVert #1 \rVert}
\newcommand{\normF}[1]{\lVert #1 \rVert_F}
\newcommand{\tr}{\operatorname{tr}}
\newcommand{\rank}{\operatorname{rank}}
\newcommand{\spanof}{\operatorname{span}}

\title{\vspace{-1.2em}The Playlist Operator\\
\large a matrix that plays music, whose picture is an image of your choice}
\author{}
\date{}

\begin{document}
\maketitle
\vspace{-3.5em}

\section{What we want to be true}

Take $N$ songs, each a vector of audio samples
$\vx_1,\dots,\vx_N \in \R^D$ (in the current build $N=6$ and
$D \approx 5.5\times10^5$ at $1$\,kHz). We want one matrix
$\mA \in \R^{D\times D}$ that \emph{is} the playlist:
\begin{equation}
  \mA\,\vx_n = \vx_{n+1} \quad (n=1,\dots,N-1),
  \qquad
  \mA\,\vx_N = \vx_1 .
  \label{eq:playlist}
\end{equation}
Seed it with any song and repeated multiplication plays the playlist,
forever, in a cycle.

Collect the songs as the columns of
$\mX = [\,\vx_1\ \vx_2\ \cdots\ \vx_N\,] \in \R^{D\times N}$, and let
$\mS \in \R^{N\times N}$ be the \emph{cyclic shift} permutation, e.g.\ for
$N=4$:
\begin{equation}
  \mS =
  \begin{pmatrix}
    0 & 0 & 0 & 1\\
    1 & 0 & 0 & 0\\
    0 & 1 & 0 & 0\\
    0 & 0 & 1 & 0
  \end{pmatrix},
  \qquad
  \mX\mS = [\,\vx_2\ \vx_3\ \vx_4\ \vx_1\,].
\end{equation}
(Column $n$ of $\mX\mS$ is the successor of song $n$; the top-right $1$
closes the loop.) All $N$ statements in \eqref{eq:playlist} then compress
into one matrix equation:
\begin{equation}
  \boxed{\ \mA\mX = \mX\mS\ }
  \label{eq:AX}
\end{equation}

\paragraph{The key count.}
Equation~\eqref{eq:AX} imposes $DN$ scalar constraints on the $D^2$ unknown
entries of $\mA$, and $N \ll D$. The songs span only an $N$-dimensional
subspace $\mathcal S = \spanof(\vx_1,\dots,\vx_N) \subset \R^D$, so
\eqref{eq:AX} dictates how $\mA$ acts on $\mathcal S$ and says
\emph{nothing} about the orthogonal complement $\mathcal S^{\perp}$, which
has dimension $D-N$. Roughly $D(D-N) \approx 3\times10^{11}$ degrees of
freedom remain. That freedom is where the artwork will live.

\section{Solving \texorpdfstring{$\mA\mX = \mX\mS$}{AX = XS}}

\subsection{Warm-up: projecting onto song space, and why
\texorpdfstring{$\mG=\mX^\top\mX$}{G = XtX} appears}

Given an arbitrary vector $\vv \in \R^D$, which combination of songs
$\mX\vc$ (coefficients $\vc \in \R^N$) is closest to it? Minimize the
squared distance, expanding it fully:
\begin{align}
  J(\vc)
  &= \norm{\mX\vc - \vv}^2
   = (\mX\vc - \vv)^\top(\mX\vc - \vv) \notag\\
  &= \vc^\top\mX^\top\mX\,\vc \;-\; \vc^\top \mX^\top\vv
     \;-\; \vv^\top\mX\,\vc \;+\; \vv^\top\vv \notag\\
  &= \vc^\top \mG\, \vc \;-\; 2\,\vc^\top \mX^\top \vv \;+\; \norm{\vv}^2,
  \qquad \mG \equiv \mX^\top\mX \in \R^{N\times N},
\end{align}
using that $\vv^\top\mX\vc$ is a scalar, hence equal to its own transpose
$\vc^\top\mX^\top\vv$. Setting the gradient to zero,
\begin{equation}
  \nabla_{\vc} J = 2\mG\vc - 2\mX^\top\vv = \bm{0}
  \quad\Longrightarrow\quad
  \mG\,\vc = \mX^\top\vv
  \quad\Longrightarrow\quad
  \vc^\star = \mG^{-1}\mX^\top\vv .
  \label{eq:normaleq}
\end{equation}
$\mG$ is the \emph{Gram matrix} of the playlist: $G_{ij} = \vx_i^\top\vx_j$
holds every pairwise song correlation. It is invertible iff the songs are
linearly independent (for real recordings, always; measured condition
number here: $1.01$ --- long audio vectors are nearly orthogonal). The
matrix $\mX^\top\vv$ \emph{correlates} the input with each song; $\mG^{-1}$
\emph{unmixes} the overlap between songs.

Reassembling the best approximation $\mX\vc^\star$ defines the orthogonal
projector onto song space and its complement:
\begin{equation}
  \mP = \mX\mG^{-1}\mX^\top, \qquad
  \mPp = \mI - \mX\mG^{-1}\mX^\top,
  \qquad
  \mPp\,\vx_n = \bm{0} \ \ \text{for every song.}
  \label{eq:projectors}
\end{equation}
($\mPp\mX = \mX - \mX\mG^{-1}\mG = \mX - \mX = \bm{0}$.)

\subsection{The minimum-norm solution, by setting a gradient to zero}

Among all solutions of \eqref{eq:AX} we single out the smallest one ---
minimum Frobenius norm ($\normF{\mA}^2 = \sum_{ij}A_{ij}^2$):
\begin{equation}
  \mAo = \arg\min_{\mA}\ \tfrac12\normF{\mA}^2
  \quad\text{subject to}\quad \mA\mX = \mX\mS .
\end{equation}
Introduce a Lagrange multiplier matrix $\mLam \in \R^{D\times N}$ (one
multiplier per scalar constraint) and form the Lagrangian
\begin{equation}
  \mathcal{L}(\mA,\mLam)
  = \tfrac12\normF{\mA}^2 \;-\; \tr\!\bigl(\mLam^\top(\mA\mX - \mX\mS)\bigr).
\end{equation}
Set the gradient with respect to $\mA$ to zero (using
$\nabla_{\mA}\,\tfrac12\normF{\mA}^2 = \mA$ and
$\nabla_{\mA}\tr(\mLam^\top\mA\mX) = \mLam\mX^\top$):
\begin{equation}
  \nabla_{\mA}\mathcal{L} = \mA - \mLam\mX^\top = \bm{0}
  \quad\Longrightarrow\quad
  \mA = \mLam\mX^\top .
  \label{eq:stationarity}
\end{equation}
So the optimum necessarily has rank $\le N$ and row space inside
$\spanof(\mX)$ --- the constraint only ``pays'' for structure aligned with
the songs; any component orthogonal to them would add norm without helping.
Substitute \eqref{eq:stationarity} into the constraint to find $\mLam$:
\begin{equation}
  \mLam\underbrace{\mX^\top\mX}_{\mG} = \mX\mS
  \quad\Longrightarrow\quad
  \mLam = \mX\mS\,\mG^{-1}
  \quad\Longrightarrow\quad
  \boxed{\ \mAo = \mX\mS\,\mG^{-1}\mX^\top\ }
  \label{eq:A0}
\end{equation}
Check it, step by step:
\begin{equation}
  \mAo\mX
  = \mX\mS\,\mG^{-1}\underbrace{\mX^\top\mX}_{=\ \mG}
  = \mX\mS\,\underbrace{\mG^{-1}\mG}_{=\ \mI_N}
  = \mX\mS. \qquad\checkmark
\end{equation}
Read $\mAo\vv = \mX\mS\,(\mG^{-1}\mX^\top\vv)$ right to left: correlate the
input with every song ($\mX^\top\vv$), unmix ($\mG^{-1}$, exactly the
least-squares coefficients \eqref{eq:normaleq}), advance one playlist step
($\mS$), rebuild audio ($\mX$).

\subsection{The full solution family}

If $\mA$ and $\mA'$ both solve \eqref{eq:AX}, their difference kills every
song: $(\mA-\mA')\mX = \bm{0}$. A matrix annihilates $\mathcal S$ exactly
when it factors through the projector $\mPp$. Hence the complete solution
set is the affine family
\begin{equation}
  \boxed{\ \mA = \mAo + \mZ\,\mPp,
  \qquad \mZ \in \R^{D\times D} \ \text{arbitrary.}\ }
  \label{eq:family}
\end{equation}
The free part is \emph{inaudible by construction}:
$\mZ\mPp\,\vx_n = \mZ\,\bm 0 = \bm 0$, so playback is exact for every
choice of $\mZ$ --- the code measures
$\max_n \norm{\mA\vx_n - \vx_{n+1}}/\norm{\vx_{n+1}} \approx 10^{-14}$,
i.e.\ floating-point roundoff, always.

\section{From matrix to picture: blocks, pooling, and two honest views}

$\mA$ has $D^2 \approx 3\times10^{11}$ entries; a print has
$P\times P \approx 10^6$ pixels. Each pixel must summarize a $k\times k$
\emph{block} of adjacent matrix entries, $k = \lfloor D/P\rfloor$ --- block
$(p,q)$ means rows $pk,\dots,(p{+}1)k-1$ and columns $qk,\dots,(q{+}1)k-1$
of $\mA$, i.e.\ ``what the operator does to time-window $q$ of the input,
as heard in time-window $p$ of the output.''

Averaging over blocks is one matrix operation. Define the \emph{pooling
matrix} $\pool \in \R^{P\times D}$: row $p$ contains $1/k$ over block $p$'s
samples and $0$ elsewhere. A toy example with $D=6$, $P=2$, $k=3$:
\begin{equation}
  \pool = \begin{pmatrix}
   \tfrac13 & \tfrac13 & \tfrac13 & 0 & 0 & 0\\[2pt]
   0 & 0 & 0 & \tfrac13 & \tfrac13 & \tfrac13
  \end{pmatrix},
\end{equation}
so $\pool\vv$ averages consecutive triples of $\vv$. $\pool$ is fully
determined by two numbers: the signal length $D$ and the chosen print
resolution $P$. Pooling rows \emph{and} columns of $\mA$ gives the
$P \times P$ \textbf{block-mean picture}
\begin{equation}
  \mM = \pool\,\mA\,\pool^\top,
  \qquad
  M_{pq} = \frac{1}{k^2}\!\!\sum_{i \in \text{block } p}\
  \sum_{j \in \text{block } q}\!\! A_{ij}.
  \label{eq:mean}
\end{equation}
Every pixel is \emph{exactly} the average of its block of true matrix
entries --- the canonical downsampling of an image too large to print, not
an approximation with error. A second, equally exact per-block statistic is
the \textbf{block energy} (root-mean-square),
$E_{pq} = \frac1{k^2}\sum_{i\in p, j\in q} A_{ij}^2$.

\emph{Why ever look at two views of one matrix?} Audio oscillates around
zero, so within any block the entries of $\mAo$ nearly cancel: the mean
picture of $\mAo$ is almost blank, while its squares do not cancel --- the
energy picture is where the music's own block structure (loudness
correlations, silence bands) is visible. Both are honest; they answer
different questions. The printed artwork uses the mean channel
\eqref{eq:mean}, because it is \emph{linear} in $\mA$ --- and that
linearity is what the next section exploits. Neither is ever computed from
the dense $\mA$: with $\mA = \bm{L}\bm{R}^\top$ kept factored,
$\mM = (\pool\bm{L})(\pool\bm{R})^\top$ and $E_{pq}$ streams through small
per-block Gram matrices.

\section{Choosing the picture}

Substitute the family \eqref{eq:family}, with a low-rank free part
$\mZ = \mU\mV^\top$ ($\mU,\mV \in \R^{D\times q}$), into the mean picture
\eqref{eq:mean}; linearity separates music from freedom:
\begin{equation}
  \mM
  = \underbrace{\pool\,\mAo\,\pool^\top}_{\ \mM_0 \text{ (fixed by the music, nearly blank)}}
  \;+\; \bigl(\pool\,\mU\bigr)\bigl(\pool\,\mPp\mV\bigr)^{\!\top}.
  \label{eq:Mdecomp}
\end{equation}
A target image $\mT$ has $P^2 \approx 10^6$ numbers; $\mZ$ has
$\approx 3\times10^{11}$ inaudible degrees of freedom. The counting is
overwhelming: (almost) \emph{any} picture is reachable. We want
\begin{equation}
  (\pool\mU)(\pool\,\mPp\mV)^{\!\top}
  \;=\; \mT_{\!\text{amp}} - \mM_0 \;\equiv\; \mDel,
  \label{eq:goal}
\end{equation}
where $\mT_{\!\text{amp}}$ is the target after centering and amplitude
scaling (a display is affine-invariant, so nothing is lost). Three steps:

\paragraph{1. Compress the image (SVD).} The truncated singular value
decomposition
\begin{equation}
  \mDel \approx \mW_{\!q}\mSig_q\mY_{\!q}^{\top}
  = \bigl(\mW_{\!q}\mSig_q^{1/2}\bigr)
    \bigl(\mY_{\!q}\mSig_q^{1/2}\bigr)^{\!\top}
  \equiv \widetilde{\mU}\,\widetilde{\mV}^{\top},
  \qquad \widetilde{\mU}, \widetilde{\mV} \in \R^{P\times q},
\end{equation}
is, by the Eckart--Young theorem, the best rank-$q$ approximation of the
residual image. Its role is purely \emph{compression} --- writing the
picture as $q$ outer products so that $\mZ$ stays storable ($D\times q$
instead of $D\times D$). Smooth cloud images compress to small $q$
($\sim\!150$, error $<1\%$); images full of sharp lines need more.

\paragraph{2. Lift pixels to samples.} We need sample-level columns with
prescribed pooled values: $\pool\mU = \widetilde{\mU}$. Any \emph{right
inverse} $\liftop \in \R^{D\times P}$ of the pooling matrix
($\pool\liftop = \mI_P$) does this. The simplest is ``hold each pixel value
constant across its $k$ samples'' ($\liftop = k\,\pool^\top$; averaging a
constant returns it). We use a smooth variant --- linear interpolation
between block centers plus a piecewise-constant correction --- which still
satisfies $\pool\liftop = \mI_P$ exactly but avoids blocky columns. Set
$\mU = \liftop\widetilde{\mU}$, $\mV = \liftop\widetilde{\mV}$.

\paragraph{3. Absorb the projector (a fixed-point iteration).} The picture
uses $\pool\mPp\mV$, not $\pool\mV$:
\begin{equation}
  \pool\,\mPp\mV
  = \pool\mV - \pool\,\mX\mG^{-1}\mX^{\top}\mV
  = \widetilde{\mV} - \underbrace{\pool\,\mX\mG^{-1}\mX^\top\mV}_{\text{rank}\ \le\, N,\ \text{small}} .
\end{equation}
The disturbance is the pooled shadow of a rank-$N$ matrix, and it is small
because generic image columns are nearly orthogonal to the $N$ specific
audio signals. Remove it by iterating
\begin{equation}
  \mV \ \leftarrow\ \mV + \liftop\bigl(\widetilde{\mV} - \pool\,\mPp\mV\bigr).
  \label{eq:fixpoint}
\end{equation}
Writing the pooled error $\ve = \widetilde{\mV} - \pool\mPp\mV$, one step
of \eqref{eq:fixpoint} maps
$\ve \leftarrow \pool\,\mX\mG^{-1}\mX^\top\liftop\,\ve$: the error is fed
through an operator of rank $\le N$ and norm well below one, so it shrinks
geometrically; two or three iterations reach sub-percent accuracy. (This is
\emph{iterative refinement} --- driving a residual to zero, as in Richardson
iteration for linear systems --- not power iteration, which does the
opposite: amplifies a dominant eigenvector.)

Because $\mV$ enters the operator only through $\mPp\mV$, playback remains
exact at every iterate; measured: image error $<1\%$, playback error
$\sim\!10^{-14}$, simultaneously.

\section*{The artwork in one line}

\begin{equation*}
  \mA \;=\;
  \underbrace{\mX\mS\,\mG^{-1}\mX^\top}_{\text{plays the playlist}}
  \;+\;
  \underbrace{\mU\bigl(\mPp\mV\bigr)^{\!\top}}_{\text{draws the picture}},
  \qquad
  \mA\vx_n = \vx_{n+1}\ \text{exactly},
  \qquad
  \pool\mA\pool^{\!\top} \approx \mT .
\end{equation*}

\end{document}
